% !TEX root = ../main.tex
\documentclass[../main.tex]{subfiles}
\graphicspath{{\subfix{../images/}}}

\begin{document}

\subsection{Жанр bullet hell: особенности и типовые механики}

Жанр компьютерных игр <<bullet hell>> характеризуется большим количеством снарядов на экране, сложными и часто симметричными паттернами их движения, а также высокими требованиями к реакции игрока. В отличие от классических аркадных шутеров, где основное внимание уделяется поражению врагов, в bullet hell важнейшим навыком становится уклонение от большого числа одновременно движущихся снарядов.

Типовые механики, встречающиеся в играх этого жанра:
\begin{itemize}
	\item \textbf{Плотные линейные залпы} — снаряды движутся параллельными рядами с небольшим интервалом.
	\item \textbf{Лучи и вереницы пуль} — пули выпускаются последовательно по одной траектории, создавая «нитку».
	\item \textbf{Круговые и веерные залпы} — снаряды разлетаются из центра под равными углами.
	\item \textbf{Прицельные снаряды} — пуля направляется в текущее положение игрока на момент выстрела.
	\item \textbf{Бомбы (медленные крупные снаряды)} — после достижения определённой точки взрываются, порождая новые снаряды.
	\item \textbf{Сложные паттерны} — комбинации перечисленных элементов, синхронизированные во времени.
\end{itemize}

Реализация таких механик требует от разработчика эффективного управления динамическими массивами (постоянное добавление и удаление снарядов), точных и быстрых алгоритмов проверки столкновений, а также удобного способа описания последовательностей атак.

\subsection{Обзор существующих игровых движков и библиотек для 2D-игр}

Для создания 2D-игр на C++ существует несколько популярных библиотек и фреймворков. Наиболее распространённые из них представлены в таблице~\ref{tab:game_engines}.

\begin{table}[H]
	\centering
	\caption{Сравнение библиотек для 2D-игр на C++}
	\label{tab:game_engines}
	\begin{tabularx}{\textwidth}{|>{\raggedright\arraybackslash}p{2.2cm}|>{\raggedright\arraybackslash}p{2.2cm}|X|X|>{\raggedright\arraybackslash}p{2.5cm}|}
		\hline
		\textbf{Биб\-ли\-о\-те\-ка} & \textbf{Лицензия} & \textbf{Аппаратное ускорение} & \textbf{Кросс\-плат\-фор\-мен\-ность} & \textbf{Сложность освоения} \\
		\hline
		SFML 2.x & zlib & Да (OpenGL) & Win / Linux / Mac & Низкая \\
		\hline
		SDL2 & zlib & Да (OpenGL / Direct3D) & Win / Linux / Mac / Android / iOS & Средняя \\
		\hline
		Raylib & zlib & Да (OpenGL) & Win / Linux / Mac / Android / Web & Низкая \\
		\hline
		GLFW + OpenGL & zlib/libpng & Да (OpenGL) & Win / Linux / Mac & Высокая \\
		\hline
	\end{tabularx}
\end{table}

\textbf{SFML} предоставляет удобные классы для графики, звука, сети и ввода, но требует установки дополнительных библиотек и может быть избыточным для простых проектов. \textbf{SDL2} — низкоуровневая библиотека, дающая полный контроль, однако её использование сопряжено с большим объёмом кода для инициализации и управления ресурсами. \textbf{Raylib} позиционируется как библиотека для обучения и прототипирования: она сочетает простоту использования, хорошую документацию и достаточную производительность. В отличие от GLFW (которая лишь создаёт окно и контекст OpenGL), Raylib предоставляет готовые функции для рисования примитивов, работы с текстурами и шрифтами, а также встроенные функции для работы со временем и векторами.

\subsection{Выбор языка программирования и библиотеки}

Для разработки был выбран язык \textbf{C++20} по следующим причинам:
\begin{itemize}
	\item Высокая производительность и контроль над памятью, необходимые для обработки большого количества снарядов (до нескольких сотен одновременно).
	\item Поддержка современного стандарта C++20, включающего \texttt{std::format} для удобного форматирования строк интерфейса и \texttt{std::span}, \texttt{std::ranges} для работы с контейнерами.
	\item Кроссплатформенность и наличие качественных компиляторов (GCC, Clang, MSVC).
\end{itemize}

Библиотекой для графики, ввода и оконного управления выбрана \textbf{Raylib} (версия 4.0 и выше). Обоснование выбора:
\begin{itemize}
	\item Лицензия zlib/libpng, позволяющая свободное использование в любых проектах.
	\item Отсутствие внешних зависимостей (кроме системных библиотек OpenGL, GLFW, ALSA и т.п.).
	\item Простой и интуитивно понятный API: для отрисовки спрайта достаточно одной функции \texttt{DrawTexturePro()}, а для обработки клавиш — \texttt{IsKeyDown()}.
	\item Встроенная поддержка работы со временем (\texttt{GetFrameTime()}) и математических функций для работы с векторами.
	\item Активное сообщество и хорошая документация.
\end{itemize}

\subsection{Требования к разрабатываемой игре}

На основе анализа жанра и возможностей инструментов были сформулированы следующие функциональные и нефункциональные требования.

\textbf{Функциональные требования:}
\begin{enumerate}
	\item Игра должна содержать меню выбора уровня с возможностью переключения между доступными уровнями.
	\item Игрок управляет персонажем с помощью клавиш WASD или стрелок; скорость движения нормализована для равномерного перемещения по диагонали.
	\item Игрок имеет показатель здоровья (HP); при столкновении со снарядом HP уменьшается, а игрок становится неуязвимым на короткое время (0,5 с).
	\item Снаряды представляют собой круги, движущиеся с постоянной скоростью по заданной траектории. При выходе за границы игрового мира снаряды удаляются.
	\item Уровни описываются через временные события (атаки), которые порождают снаряды. События могут создавать новые события (например, взрыв бомбы).
	\item После завершения всех событий и исчезновения всех снарядов уровень считается пройденным, игра переходит в состояние \texttt{LEVEL\_COMPLETED}.
	\item При падении здоровья до 0 игрок проигрывает, состояние игры — \texttt{GAME\_OVER}.
	% \item Должен быть предусмотрен режим отладки (включается клавишей F3), отображающий дополнительную информацию (FPS, точное положение игрока, хитбоксы, количество коллизий).
\end{enumerate}

\textbf{Нефункциональные требования:}
\begin{itemize}
	\item Частота кадров должна быть не менее 60 FPS при 300 активных снарядах на современном оборудовании.
	\item Исходный код должен быть структурирован, разделён на логические модули (game.cpp, render.cpp, main.cpp).
	\item Сборка должна осуществляться с помощью системы CMake или Premake, поддерживаться в среде Linux.
	\item Интерфейс игры должен быть адаптивным к изменению размера окна.
\end{itemize}

\subsection{Постановка задачи}

На основе сформулированных требований необходимо разработать программную систему, реализующую игру в жанре bullet hell. В рамках курсовой работы требуется:

\begin{enumerate}
	\item Спроектировать объектную модель, включающую классы \texttt{Player}, \texttt{Projectile}, \texttt{AttackEvent} и перечисление состояний игры (\texttt{GameState}).
	\item Реализовать игровой цикл, поддерживающий переходы между состояниями \texttt{MENU}, \texttt{PLAYING}, \texttt{LEVEL\_COMPLETED}, \texttt{GAME\_OVER}.
	\item Разработать механику движения игрока с нормализацией скорости и ограничением границами мира.
	\item Реализовать систему событий, позволяющую задавать атаки в виде функций с привязкой к абсолютному времени уровня.
	\item Создать библиотеку вспомогательных функций для часто используемых паттернов атак: \texttt{add\_beam\_events}, \texttt{spawn\_circular\_burst}, \texttt{spawn\_bullet\_line}, \texttt{add\_fan\_attack}.
	\item Разработать два демонстрационных уровня: <<Beginning>> (обучающий) и <<Way>> (с более сложными паттернами).
	\item Реализовать отрисовку с преобразованием мировых координат в экранные (с сохранением пропорций и центрированием).
	\item Подготовить инструкцию по сборке и запуску на платформе Linux.
\end{enumerate}

Решение перечисленных задач позволит получить работающий прототип игры, демонстрирующий основные принципы организации игровой логики на C++ с использованием Raylib.

\end{document}
